% GENERATED FILE. Edit canonical lessons, then run npm run generate:books.
% canonical-source-hash: fnv1a64:ead77e270f4ab156
% canonical-lessons: PT-C329-antepassado, PT-C329-batalha, PT-C329-invadir, PT-C329-conquistar, PT-C329-revolucao

\chapter{Páginas da história}
\label{ch:pt-paginas-da-historia}
\hlchaptermodality{329}

\begin{chapteropening}
\textbf{By the end of this chapter:} \emph{I can tell history: ancestors, battles, invasions, conquests and revolutions.}

Complete the last lesson of chapter 329: I can tell history: ancestors, battles, invasions, conquests and revolutions.
\end{chapteropening}

\section[o antepassado]{o antepassado — the ancestor}

\label{lesson:PT-C329-antepassado}



\begin{quote}
Before the new one: say the Portuguese for the shipwreck, then the Portuguese for the treasure.
\end{quote}

\subsection*{You'll want to know: o antepassado}
\textbf{o antepassado} — \textquotedblleft{}the ancestor\textquotedblright{}.

\begin{quote}
Um antepassado meu viveu nesta casa.
\end{quote}

\begin{quote}
\emph{An ancestor of mine lived in this house.}
\end{quote}

Usually plural: \emph{os nossos antepassados}.

\begin{cousinweb}
\emph{antepassado} joins \emph{ante} \textquotedblleft{}before\textquotedblright{}, from Latin \emph{ante} (Lewis \& Short), and \emph{passado} \textquotedblleft{}past\textquotedblright{}: one who lived before us.
\end{cousinweb}

\begin{grammarlens}[title={What you've built}]
One more word for telling, explaining and describing at B1.
\end{grammarlens}

\subsection*{Your turn}
\begin{itemize}
\raggedright
  \item \emph{Say it:} \emph{o antepassado}
  \item \emph{Say it:} \emph{o antepassado}, once more
  \item \emph{Say it:} a sentence of your own with \emph{o antepassado}
  \item \emph{Recall:} say the Portuguese for the shipwreck, then the Portuguese for the treasure, then say \emph{o antepassado} again
\end{itemize}

\begin{tcolorbox}[breakable,colback=teal!4,colframe=teal!35!black,title={Before you move on}]
What does o antepassado mean? (\textquotedblleft{}The ancestor\textquotedblright{}.) Say it once more.
\end{tcolorbox}
\section[a batalha]{a batalha — the battle}

\label{lesson:PT-C329-batalha}



\begin{quote}
Before the new one: say the Portuguese for the treasure, then the Portuguese for the ancestor.
\end{quote}

\subsection*{You'll want to know: a batalha}
\textbf{a batalha} — \textquotedblleft{}the battle\textquotedblright{}.

\begin{quote}
A batalha durou um dia inteiro.
\end{quote}

\begin{quote}
\emph{The battle lasted a whole day.}
\end{quote}

Also any hard struggle: \emph{uma batalha contra a doença}.

\begin{cousinweb}
\emph{batalha} comes from Late Latin \emph{battālia} \textquotedblleft{}fighting, fencing\textquotedblright{} (EtymDB), built on \emph{battuere} \textquotedblleft{}to beat\textquotedblright{}.
\end{cousinweb}

\begin{grammarlens}[title={What you've built}]
One more word for telling, explaining and describing at B1.
\end{grammarlens}

\subsection*{Your turn}
\begin{itemize}
\raggedright
  \item \emph{Say it:} \emph{a batalha}
  \item \emph{Say it:} \emph{a batalha}, once more
  \item \emph{Say it:} a sentence of your own with \emph{a batalha}
  \item \emph{Recall:} say the Portuguese for the treasure, then the Portuguese for the ancestor, then say \emph{a batalha} again
\end{itemize}

\begin{tcolorbox}[breakable,colback=teal!4,colframe=teal!35!black,title={Before you move on}]
What does a batalha mean? (\textquotedblleft{}The battle\textquotedblright{}.) Say it once more.
\end{tcolorbox}
\section[invadir]{invadir — to invade}

\label{lesson:PT-C329-invadir}



\begin{quote}
Before the new one: say the Portuguese for the ancestor, then the Portuguese for the battle.
\end{quote}

\subsection*{You'll want to know: invadir}
\textbf{invadir} — \textquotedblleft{}to invade\textquotedblright{}.

\begin{quote}
O exército invadiu o país vizinho.
\end{quote}

\begin{quote}
\emph{The army invaded the neighbouring country.}
\end{quote}

Also of crowds: \emph{os adeptos invadiram o campo}.

\begin{cousinweb}
\emph{invadir} was borrowed from Latin \emph{invādere} \textquotedblleft{}to go into, to attack\textquotedblright{} (EtymDB), built on \emph{vādere} \textquotedblleft{}to go\textquotedblright{}.
\end{cousinweb}

\begin{grammarlens}[title={What you've built}]
One more word for telling, explaining and describing at B1.
\end{grammarlens}

\subsection*{Your turn}
\begin{itemize}
\raggedright
  \item \emph{Say it:} \emph{invadir}
  \item \emph{Say it:} \emph{invadir}, once more
  \item \emph{Say it:} a sentence of your own with \emph{invadir}
  \item \emph{Recall:} say the Portuguese for the ancestor, then the Portuguese for the battle, then say \emph{invadir} again
\end{itemize}

\begin{tcolorbox}[breakable,colback=teal!4,colframe=teal!35!black,title={Before you move on}]
What does invadir mean? (\textquotedblleft{}To invade\textquotedblright{}.) Say it once more.
\end{tcolorbox}
\section[conquistar]{conquistar — to conquer; to win (over)}

\label{lesson:PT-C329-conquistar}



\begin{quote}
Before the new one: say the Portuguese for the battle, then the Portuguese for to invade.
\end{quote}

\subsection*{You'll want to know: conquistar}
\textbf{conquistar} — \textquotedblleft{}to conquer; to win (over)\textquotedblright{}.

\begin{quote}
O rei conquistou a cidade.
\end{quote}

\begin{quote}
\emph{The king conquered the city.}
\end{quote}

Also to win someone over: \emph{conquistou o público}.

\begin{cousinweb}
\emph{conquistar} is usually traced to Vulgar Latin \emph{conquaesītāre}, from Latin \emph{conquīrere} \textquotedblleft{}to seek out\textquotedblright{}, built on \emph{quaerere} \textquotedblleft{}to seek\textquotedblright{} (Lewis \& Short).
\end{cousinweb}

\begin{grammarlens}[title={What you've built}]
One more word for telling, explaining and describing at B1.
\end{grammarlens}

\subsection*{Your turn}
\begin{itemize}
\raggedright
  \item \emph{Say it:} \emph{conquistar}
  \item \emph{Say it:} \emph{conquistar}, once more
  \item \emph{Say it:} a sentence of your own with \emph{conquistar}
  \item \emph{Recall:} say the Portuguese for the battle, then the Portuguese for to invade, then say \emph{conquistar} again
\end{itemize}

\begin{tcolorbox}[breakable,colback=teal!4,colframe=teal!35!black,title={Before you move on}]
What does conquistar mean? (\textquotedblleft{}To conquer; to win (over)\textquotedblright{}.) Say it once more.
\end{tcolorbox}
\section[a revolução]{a revolução — the revolution}

\label{lesson:PT-C329-revolucao}



\begin{quote}
Before the new one: say the Portuguese for to invade, then the Portuguese for to conquer.
\end{quote}

\subsection*{You'll want to know: a revolução}
\textbf{a revolução} — \textquotedblleft{}the revolution\textquotedblright{}.

\begin{quote}
A revolução aconteceu sem violência.
\end{quote}

\begin{quote}
\emph{The revolution happened without violence.}
\end{quote}

Also a turn of a wheel or a planet.

\begin{cousinweb}
\emph{revolução} was borrowed from Late Latin \emph{revolūtiō} \textquotedblleft{}a turning back\textquotedblright{} (EtymDB), from \emph{revolvere} \textquotedblleft{}to roll back\textquotedblright{}, built on \emph{volvere} \textquotedblleft{}to roll\textquotedblright{}.
\end{cousinweb}

\begin{grammarlens}[title={What you've built}]
One more word for telling, explaining and describing at B1.
\end{grammarlens}

\subsection*{Your turn}
\begin{itemize}
\raggedright
  \item \emph{Say it:} \emph{a revolução}
  \item \emph{Say it:} \emph{a revolução}, once more
  \item \emph{Say it:} a sentence of your own with \emph{a revolução}
  \item \emph{Recall:} say the Portuguese for to invade, then the Portuguese for to conquer, then say \emph{a revolução} again
\end{itemize}

\begin{tcolorbox}[breakable,colback=teal!4,colframe=teal!35!black,title={Before you move on}]
What does a revolução mean? (\textquotedblleft{}The revolution\textquotedblright{}.) Say it once more.
\end{tcolorbox}
